\chapter{Bariatric Surgery}

\section*{Definition and Overview}
Bariatric surgery, also called weight-loss surgery or metabolic surgery, is a group of operations that help people with severe obesity lose a large amount of weight. The procedures work by reducing the size of the stomach, changing how food passes through the gut, and altering the gut hormones that control hunger and blood sugar.

The main operations are the gastric sleeve (sleeve gastrectomy), which removes most of the stomach to leave a narrow tube, and the gastric bypass, which creates a small stomach pouch and re-routes part of the small intestine. The adjustable gastric band, which tightens the stomach with an inflatable band, is less commonly used now. Bariatric surgery is the most effective available treatment for severe obesity and can improve or resolve related conditions such as type 2 diabetes. It is a major operation, requires lifelong follow-up, and is offered only after careful assessment.

\section*{Epidemiology}
Severe obesity (a body mass index of 40 or above, or 35 with complications) affects a significant minority of Australian adults. Bariatric surgery is increasingly common, with many thousands of procedures performed in Australia each year, most of them gastric sleeves. The typical patient is in their 40s or 50s, and more women than men undergo the surgery. Demand has grown because severe obesity is linked with diabetes, heart disease, joint problems, and reduced life expectancy, and because surgery produces far greater weight loss than diet and exercise alone in this group.

\section*{Aetiology and Pathophysiology}
Severe obesity is a chronic disease driven by a mix of genetics, behaviour, environment, and hormones that regulate appetite. Bariatric surgery treats it by several mechanisms:

\begin{itemize}
    \item \textbf{Restriction:} a smaller stomach holds less food, so intake is limited
    \item \textbf{Malabsorption (bypass procedures):} food bypasses part of the small intestine, reducing absorption of some energy and nutrients
    \item \textbf{Hormonal change:} altered gut hormones lower appetite, increase fullness, and improve insulin handling -- often before much weight is lost
    \item \textbf{Indications:} surgery is considered for a BMI of 40 or above, or 35 or above with serious obesity-related complications such as type 2 diabetes, sleep apnoea, or severe joint disease, when structured weight-loss attempts have not worked
\end{itemize}

\section*{Clinical Features}
The features are those of severe obesity and its complications, which bariatric surgery aims to treat:

\begin{itemize}
    \item Type 2 diabetes and abnormal blood sugar
    \item High blood pressure and abnormal cholesterol
    \item Obstructive sleep apnoea, with loud snoring and daytime sleepiness
    \item Joint pain and osteoarthritis, especially of the knees and hips
    \item Acid reflux and heartburn
    \item Fatty liver disease
    \item Reduced mobility, breathlessness on exertion, and fatigue
    \item Psychological effects, including low mood and weight stigma
\end{itemize}

\section*{Differential Diagnosis}
Before surgery, the assessment distinguishes uncomplicated severe obesity from problems that need separate treatment or change the surgical plan:

\begin{itemize}
    \item Obesity with versus without obesity-related complications, which influence urgency and benefit
    \item Genetic or hormonal causes of weight gain, such as an underactive thyroid
    \item Medication-induced weight gain (for example, some antidepressants and corticosteroids)
    \item Binge eating or other disordered eating, which may need treatment before or instead of surgery
    \item Medical contraindications to surgery, such as uncontrolled mental illness, active substance misuse, or major organ failure
\end{itemize}

\section*{When to Seek Medical Care}
Bariatric surgery is a major procedure, and careful follow-up after leaving hospital is essential.

\textbf{Call 000 (or local emergency services) or contact your surgical team urgently for:}
\begin{itemize}
    \item Severe or worsening abdominal pain, especially in the upper abdomen or chest
    \item Persistent vomiting or inability to keep fluids down
    \item Fever, chills, or shaking
    \item Difficulty swallowing, or new and severe nausea after eating
    \item A fast heart rate, shortness of breath, or light-headedness
    \item Wound redness, discharge, or separation at the incision sites
    \item Black stools, or vomiting blood
\end{itemize}

See a doctor sooner rather than later for any new symptom you are unsure about in the first weeks after surgery.

\section*{Diagnosis and Investigations}
\begin{itemize}
    \item \textbf{Weight and health assessment:} BMI, waist circumference, and a review of obesity-related complications
    \item \textbf{Blood tests:} blood glucose or HbA1c, cholesterol, thyroid, liver, and kidney function
    \item \textbf{Sleep study:} to diagnose obstructive sleep apnoea, which is common
    \item \textbf{Heart assessment:} ECG and, if indicated, further cardiac tests
    \item \textbf{Gastroscopy:} to check the stomach and oesophagus, and to treat any reflux before surgery
    \item \textbf{Dietitian and psychological assessment:} to review eating habits, readiness, and support for the lifestyle changes required
    \item \textbf{Multidisciplinary review:} the case is discussed by a team including the surgeon, physician, dietitian, and psychologist before approval
\end{itemize}

\section*{Management}
\begin{itemize}
    \item \textbf{Pre-operative preparation:} stopping smoking, optimising diabetes and blood pressure, a pre-surgery diet, and attending education sessions
    \item \textbf{Choice of procedure:} sleeve and bypass are the most common; the choice depends on the person's body weight, conditions, and preferences
    \item \textbf{Perioperative care:} keyhole (laparoscopic) techniques, blood-clot prevention, and early mobilisation
    \item \textbf{Staged diet after surgery:} clear fluids, then liquids, then pureed and soft foods, building up over several weeks as advised
    \item \textbf{Lifelong supplementation:} vitamin and mineral supplements (especially iron, B12, calcium, and vitamin D) are needed for life
    \item \textbf{Follow-up:} regular review of weight, nutrition, and complications, with dietitian and psychological support
    \item \textbf{Lifestyle change:} healthy eating, physical activity, and behavioural strategies are essential for lasting results
\end{itemize}

\section*{Complications}
\begin{itemize}
    \item Early: bleeding, infection, blood clots (deep vein thrombosis, pulmonary embolism), and leaks at the surgical joins, which can be serious
    \item Nutritional deficiencies: iron, vitamin B12, calcium, vitamin D, and protein deficiency, requiring supplementation
    \item Dumping syndrome (more after bypass): nausea, sweating, and diarrhoea after sugary food
    \item Gallstones, which may form with rapid weight loss
    \item Strictures (narrowing) or ulcers at the surgical joins, causing difficulty eating
    \item Band problems (where a band is used): slippage, erosion, or reflux
    \item Weight regain over time, and loose, excess skin that may need further surgery
    \item A small risk of death from the surgery itself, which must be weighed against the risks of untreated severe obesity
\end{itemize}

\section*{Prognosis and Outcomes}
Bariatric surgery produces substantial, sustained weight loss: on average around a quarter to a third of starting body weight, depending on the procedure. Type 2 diabetes improves or goes into remission in many people, and sleep apnoea, blood pressure, and joint pain commonly improve. Quality of life usually increases substantially. However, results depend on lifelong follow-up, supplements, and lifestyle changes, and some weight regain is common over the long term. While surgery is the most effective treatment for severe obesity, it is a major operation with real risks, and the decision is made individually after a full assessment.

\section*{Prevention and Self-Care}
\begin{itemize}
    \item For most people, weight gain is best prevented early: a balanced diet, regular physical activity, and healthy sleep habits
    \item If you are considering bariatric surgery, first try structured, professionally supported weight-loss programs
    \item Stop smoking and optimise diabetes, blood pressure, and other conditions well before surgery
    \item Attend all pre-operative education and follow-up appointments
    \item Follow the staged diet and texture guidelines strictly after surgery
    \item Take vitamin and mineral supplements exactly as prescribed, for life
    \item Build regular physical activity back up gradually, with guidance
    \item Delay pregnancy for at least 12--24 months after surgery and plan it with your team
    \item Seek help early for any new or worrying symptom after surgery
\end{itemize}

\section*{Sources and Further Reading}
The following sources provide reliable, up-to-date information on bariatric surgery.
\begin{itemize}
    \item NHS. \textit{Weight loss surgery.} https://www.nhs.uk/conditions/
    \item Mayo Clinic. \textit{Bariatric surgery.} https://www.mayoclinic.org
    \item MSD Manual. \textit{Obesity.} https://www.msdmanuals.com
    \item NICE. \textit{Obesity: identification, assessment and management.} https://www.nice.org.uk
    \item World Health Organization. \textit{Obesity.} https://www.who.int
\end{itemize}
